% blog-title: Foundations of Learning
% blog-date: 2026-06-15
% blog-description: A starter note for Tang's Machine Learning Blog, written in LaTeX with mathematical notation, code highlighting, and reproducible structure.
% blog-lang: en

\documentclass[en,nofont,11pt]{elegantbook}

\title{Foundations of Learning}
\subtitle{A LaTeX-native note for machine learning research}
\author{Tang}
\institute{Tang's Machine Learning Blog}
\date{2026-06-15}
\version{0.1.0}
\bioinfo{Topic}{Machine Learning}

\usepackage{fontspec}
\ifdefined\setCJKmainfont
\else
  \usepackage[UTF8,scheme=plain,fontset=none]{ctex}
\fi
\usepackage{xcolor}
\usepackage{listings}
\usepackage{float}
\usepackage{tcolorbox}
\tcbuselibrary{breakable, listings, skins}
\usetikzlibrary{arrows.meta,positioning,calc,decorations.pathmorphing,decorations.pathreplacing,fit,backgrounds,patterns,shapes.geometric}

\IfFontExistsTF{Microsoft YaHei}{
  \setCJKmainfont{Microsoft YaHei}
  \setCJKsansfont{Microsoft YaHei}
  \setCJKmonofont{Microsoft YaHei UI}
}{
  \IfFontExistsTF{Noto Serif CJK SC}{
    \setCJKmainfont{Noto Serif CJK SC}
    \setCJKsansfont{Noto Sans CJK SC}
    \setCJKmonofont{Noto Sans Mono CJK SC}
  }{}
}

\definecolor{blogAccent}{HTML}{1F6F8B}
\definecolor{blogInk}{HTML}{172026}
\definecolor{blogCodeBg}{HTML}{F5F7F8}
\definecolor{blogCodeBorder}{HTML}{D6E2E8}
\definecolor{blogKeyword}{HTML}{B23A48}
\definecolor{blogString}{HTML}{287D4F}
\definecolor{blogComment}{HTML}{687783}

\lstdefinestyle{blogcode}{
  basicstyle=\ttfamily\small,
  keywordstyle=\bfseries\color{blogKeyword},
  stringstyle=\color{blogString},
  commentstyle=\itshape\color{blogComment},
  numberstyle=\scriptsize\color{blogComment},
  numbers=left,
  numbersep=8pt,
  showstringspaces=false,
  breaklines=true,
  tabsize=2,
  columns=fullflexible,
  keepspaces=true
}

\newtcblisting{codeblock}[2][]{
  listing only,
  breakable,
  enhanced,
  colback=blogCodeBg,
  colframe=blogCodeBorder,
  arc=2mm,
  boxrule=0.4pt,
  left=6mm,
  right=2mm,
  top=1mm,
  bottom=1mm,
  listing options={style=blogcode, language=#2},
  #1
}

\hypersetup{
  colorlinks=true,
  linkcolor=blogAccent,
  urlcolor=blogAccent,
  citecolor=blogAccent
}


\begin{document}
\maketitle
\frontmatter
\tableofcontents
\mainmatter

\chapter{A Minimal Research Note}

This blog is designed for machine learning writing that needs clean equations, readable code, and stable PDF rendering. Each post is a standalone \verb|.tex| file under \verb|posts/| and is compiled automatically by GitHub Actions.

\section{Optimization}

A standard regularized empirical risk objective can be written as
\[
  \min_{\theta} \frac{1}{n}\sum_{i=1}^{n}\ell(f_{\theta}(x_i), y_i)
  + \lambda\|\theta\|_2^2.
\]

\section{Code}

\begin{codeblock}{Python}
from dataclasses import dataclass

@dataclass
class Experiment:
    name: str
    learning_rate: float = 3e-4

def train(config: Experiment) -> None:
    print(f"running {config.name} at lr={config.learning_rate}")

train(Experiment("tang-ml-blog"))
\end{codeblock}

\section{Bilingual Writing}

English posts may use \verb|\documentclass[en,nofont,11pt]{elegantbook}|. Chinese posts may use \verb|\documentclass[cn,nofont,11pt]{elegantbook}|. The shared preamble loads robust CJK fonts when available, so both languages can be compiled by the same pipeline.

\end{document}
